% Fragmento público generado desde propietarios íntegros.
% Residencia: 52.3.1.
% Fuente: ../incoming/radion_monografia_20260827/manuscrito/sections/04_cincuenta_critica.tex:1-105
\noindent La coordonnée critique relie la carte de Cayley, le registre APP et la phase dodécaphasique par des applications typées.\par\medskip

\subsubsection{Deux réalisations exactes de la coordonnée angulaire \(50^\circ\)}

\paragraph{Deuxième phase du système temporel dodécaphasique}

Par \eqref{eq:rueda-dodecafasica}, la suite des centres commence
\[
20^\circ,\ 50^\circ,\ 80^\circ,\ 110^\circ,\ldots.
\]
Résoudre
\[
30m-10=50
\]
donne uniquement \(m=2\). Par conséquent, \(50^\circ\) n'est ni un arrondi de \(A\),
ni un demi-tour ni le centre d'un cercle : c'est une phase typée du calendrier
dodécaphasique. La fraction de tour est
\[
\frac{50}{360}=\frac5{36},
\]
exactement le premier terme de \(S_E/T_+\).

\paragraph{Égalisation de la coordonnée critique de Cayley}

L'involution fonctionnelle
\[
\rho\longmapsto1-\bar\rho
\]
fixe
\[
\operatorname{Fix}(\rho\mapsto1-\bar\rho)
=\{\rho:\Re\rho=1/2\}.
\]
Le nombre \(1/2\) est la coordonnée réelle autoduale de l'intervalle fonctionnel
\(0\leq\Re s\leq1\). Ce n'est pas \enquote{la moitié d'une droite} au sens
longitudinal, et il ne contient pas la hauteur spectrale \(\gamma\). Pour
\[
\rho_\gamma=\frac12+i\gamma,
\]
la transformation de Cayley peut s'écrire
\begin{equation}
\tau_\gamma
=\frac{\rho_\gamma-1}{\rho_\gamma}
=\frac{\gamma+i/2}{\gamma-i/2}
\in S^1,
\qquad
\gamma=\frac12\cot\frac{\theta_\gamma}{2}.
\label{eq:cayley-critica}
\end{equation}
Ainsi, \(1/2\) fixe la couture, \(\gamma\) détermine le caractère angulaire et
\(\Gamma_9\) conserve la profondeur de ses itérations.

\paragraph{Injection affine \(j_{100}\) du registre APP dans la carte angulaire}

La complétion cubique dispose une chambre de dix marques par coordonnée. Une
face contient \(10^2=100\) positions. Nous déclarons la carte affine
\begin{equation}
j_{100}:[0,1]\longrightarrow[0,100^\circ],
\qquad j_{100}(\sigma)=100\sigma^\circ.
\label{eq:j100}
\end{equation}
Alors
\[
j_{100}\left(\frac12\right)=50^\circ.
\]
En exigeant simultanément \(j_{100}(1/2)=\theta_m\), il vient
\[
100\cdot\frac12=30m-10
\quad\Longleftrightarrow\quad m=2.
\]

\begin{theorem}[Égalisateur critique--dodécaphasique]
Dans la carte APP de cent marques, la coordonnée fixe de l'involution critique
se publie exactement dans la deuxième phase de la roue :
\begin{equation}
\boxed{j_{100}(1/2)=\theta_2=50^\circ.}
\label{eq:igualador-50}
\end{equation}
La phase \(m=2\) est unique.
\end{theorem}

Le théorème est exact une fois \eqref{eq:j100} déclarée. Son statut est
\status{formalisation nouvelle}: le demi critique, la complétion \(10^3\), la
face de cent marques et \(\theta_2\) étaient déjà construits ; le choix de cette
face comme carte critique affine réunit pour la première fois leurs domaines.

\paragraph{Compatibilité de la valeur critique avec la carte de Cayley}

La transformation de Cayley n'envoie pas \(1/2\) sur \(50^\circ\). Si
\(\gamma=0\), \eqref{eq:cayley-critica} donne
\[
\tau_0=-1=e^{i\pi},
\]
c'est-à-dire \(180^\circ\). Les \(50^\circ\) appartiennent à la carte APP de cent
marques et au calendrier dodécaphasique ; l'angle de Cayley dépend de \(\gamma\).
Cette distinction empêche de réduire tous les zéros de la fonction zêta à une seule
phase.

\begin{falsifier}
L'identification \(1/2\leftrightarrow50^\circ\) échoue si l'on élimine la
carte \eqref{eq:j100}. L'égalité scalaire \(100(1/2)=50\) ne suffit pas
à identifier à elle seule deux objets. Le théorème dépend de la face APP
distinguée et de la roue dodécaphasique.
\end{falsifier}

